% GENERATED FILE. Edit canonical lessons, then run npm run generate:books.
% canonical-source-hash: fnv1a64:702798ddd66d50e1
% canonical-lessons: IT-C396-contanti, IT-C396-saldo, IT-C396-assegno, IT-C396-versamento, IT-C396-deposito

\chapter{In banca}
\label{ch:it-in-banca}
\hlchaptermodality{396}

\begin{chapteropening}
\textbf{By the end of this chapter:} \emph{I can handle money at the bank: cash, the balance, a cheque, a payment in and a deposit.}

Complete the last lesson of chapter 396: I can handle money at the bank: cash, the balance, a cheque, a payment in and a deposit.
\end{chapteropening}

\section[i contanti]{i contanti — the cash}

\label{lesson:IT-C396-contanti}



\begin{quote}
Before the new one: say the Italian for the rate, then the Italian for the inflation.
\end{quote}

\subsection*{You'll want to know: i contanti}
\textbf{i contanti} — \textquotedblleft{}the cash\textquotedblright{}.

\begin{quote}
Posso pagare in contanti?
\end{quote}

\begin{quote}
\emph{Can I pay in cash?}
\end{quote}

Usually plural: \emph{pagare in contanti}. The opposite is \emph{con la carta}.

\begin{cousinweb}
\emph{contanti} is the plural of \emph{contante} \textquotedblleft{}counted out\textquotedblright{}, from \emph{contare} \textquotedblleft{}to count\textquotedblright{}, inherited from Latin \emph{computāre} \textquotedblleft{}to reckon\textquotedblright{} (Lewis \& Short): money you can count in your hand.
\end{cousinweb}

\begin{grammarlens}[title={What you've built}]
One more word for telling, explaining and describing at B1.
\end{grammarlens}

\subsection*{Your turn}
\begin{itemize}
\raggedright
  \item \emph{Say it:} \emph{i contanti}
  \item \emph{Say it:} \emph{i contanti}, once more
  \item \emph{Say it:} a sentence of your own with \emph{i contanti}
  \item \emph{Recall:} say the Italian for the rate, then the Italian for the inflation, then say \emph{i contanti} again
\end{itemize}

\begin{tcolorbox}[breakable,colback=teal!4,colframe=teal!35!black,title={Before you move on}]
What does i contanti mean? (\textquotedblleft{}The cash\textquotedblright{}.) Say it once more.
\end{tcolorbox}
\section[il saldo]{il saldo — the balance (of an account); the final payment}

\label{lesson:IT-C396-saldo}



\begin{quote}
Before the new one: say the Italian for the inflation, then the Italian for the cash.
\end{quote}

\subsection*{You'll want to know: il saldo}
\textbf{il saldo} — \textquotedblleft{}the balance (of an account); the final payment\textquotedblright{}.

\begin{quote}
Controllo il saldo prima di fare acquisti.
\end{quote}

\begin{quote}
\emph{I check my balance before I go shopping.}
\end{quote}

\emph{Il saldo} is also what is left to pay: \emph{pagare il saldo}. \emph{I saldi} are the sales.

\begin{cousinweb}
\emph{saldo} is built on \emph{saldare} \textquotedblleft{}to weld, to settle an account\textquotedblright{}, from the adjective \emph{saldo} \textquotedblleft{}firm\textquotedblright{}, inherited from Latin \emph{solidus} \textquotedblleft{}solid, whole\textquotedblright{} (Lewis \& Short): the account made whole.
\end{cousinweb}

\begin{grammarlens}[title={What you've built}]
One more word for telling, explaining and describing at B1.
\end{grammarlens}

\subsection*{Your turn}
\begin{itemize}
\raggedright
  \item \emph{Say it:} \emph{il saldo}
  \item \emph{Say it:} \emph{il saldo}, once more
  \item \emph{Say it:} a sentence of your own with \emph{il saldo}
  \item \emph{Recall:} say the Italian for the inflation, then the Italian for the cash, then say \emph{il saldo} again
\end{itemize}

\begin{tcolorbox}[breakable,colback=teal!4,colframe=teal!35!black,title={Before you move on}]
What does il saldo mean? (\textquotedblleft{}The balance (of an account); the final payment\textquotedblright{}.) Say it once more.
\end{tcolorbox}
\section[l'assegno]{l'assegno — the cheque}

\label{lesson:IT-C396-assegno}



\begin{quote}
Before the new one: say the Italian for the cash, then the Italian for the balance.
\end{quote}

\subsection*{You'll want to know: l'assegno}
\textbf{l'assegno} — \textquotedblleft{}the cheque\textquotedblright{}.

\begin{quote}
Oggi pochi pagano con l'assegno.
\end{quote}

\begin{quote}
\emph{Few people pay by cheque today.}
\end{quote}

\emph{Un assegno a vuoto} — a cheque that bounces.

\begin{cousinweb}
\emph{assegno} is built on \emph{assegnare} \textquotedblleft{}to assign\textquotedblright{}, a learned borrowing from Latin \emph{assignāre} \textquotedblleft{}to allot, to mark out\textquotedblright{} (Lewis \& Short), built on \emph{signum} \textquotedblleft{}a mark, a seal\textquotedblright{}.
\end{cousinweb}

\begin{grammarlens}[title={What you've built}]
One more word for telling, explaining and describing at B1.
\end{grammarlens}

\subsection*{Your turn}
\begin{itemize}
\raggedright
  \item \emph{Say it:} \emph{l'assegno}
  \item \emph{Say it:} \emph{l'assegno}, once more
  \item \emph{Say it:} a sentence of your own with \emph{l'assegno}
  \item \emph{Recall:} say the Italian for the cash, then the Italian for the balance, then say \emph{l'assegno} again
\end{itemize}

\begin{tcolorbox}[breakable,colback=teal!4,colframe=teal!35!black,title={Before you move on}]
What does l'assegno mean? (\textquotedblleft{}The cheque\textquotedblright{}.) Say it once more.
\end{tcolorbox}
\section[il versamento]{il versamento — the payment (into an account), the deposit}

\label{lesson:IT-C396-versamento}



\begin{quote}
Before the new one: say the Italian for the balance, then the Italian for the cheque.
\end{quote}

\subsection*{You'll want to know: il versamento}
\textbf{il versamento} — \textquotedblleft{}the payment (into an account), the deposit\textquotedblright{}.

\begin{quote}
Ho fatto un versamento sul conto di mio figlio.
\end{quote}

\begin{quote}
\emph{I made a payment into my son's account.}
\end{quote}

Built on \emph{versare} \textquotedblleft{}to pour; to pay in\textquotedblright{}.

\begin{cousinweb}
\emph{versamento} is built on \emph{versare} \textquotedblleft{}to pour, to pay in\textquotedblright{}, inherited from Latin \emph{versāre} \textquotedblleft{}to turn often\textquotedblright{} (Lewis \& Short), built on \emph{vertere} \textquotedblleft{}to turn\textquotedblright{}.
\end{cousinweb}

\begin{grammarlens}[title={What you've built}]
One more word for telling, explaining and describing at B1.
\end{grammarlens}

\subsection*{Your turn}
\begin{itemize}
\raggedright
  \item \emph{Say it:} \emph{il versamento}
  \item \emph{Say it:} \emph{il versamento}, once more
  \item \emph{Say it:} a sentence of your own with \emph{il versamento}
  \item \emph{Recall:} say the Italian for the balance, then the Italian for the cheque, then say \emph{il versamento} again
\end{itemize}

\begin{tcolorbox}[breakable,colback=teal!4,colframe=teal!35!black,title={Before you move on}]
What does il versamento mean? (\textquotedblleft{}The payment (into an account), the deposit\textquotedblright{}.) Say it once more.
\end{tcolorbox}
\section[il deposito]{il deposito — the deposit; the storeroom}

\label{lesson:IT-C396-deposito}



\begin{quote}
Before the new one: say the Italian for the cheque, then the Italian for the payment.
\end{quote}

\subsection*{You'll want to know: il deposito}
\textbf{il deposito} — \textquotedblleft{}the deposit; the storeroom\textquotedblright{}.

\begin{quote}
Per l'appartamento devo lasciare un deposito di due mesi.
\end{quote}

\begin{quote}
\emph{For the flat I have to leave a two-month deposit.}
\end{quote}

\emph{Il deposito bagagli} — the left-luggage office.

\begin{cousinweb}
\emph{deposito} is a learned borrowing from Latin \emph{dēpositum} \textquotedblleft{}a thing laid down for safekeeping\textquotedblright{} (Lewis \& Short), from \emph{dēpōnere} \textquotedblleft{}to lay down\textquotedblright{}, built on \emph{pōnere} \textquotedblleft{}to put\textquotedblright{}.
\end{cousinweb}

\begin{grammarlens}[title={What you've built}]
One more word for telling, explaining and describing at B1.
\end{grammarlens}

\subsection*{Your turn}
\begin{itemize}
\raggedright
  \item \emph{Say it:} \emph{il deposito}
  \item \emph{Say it:} \emph{il deposito}, once more
  \item \emph{Say it:} a sentence of your own with \emph{il deposito}
  \item \emph{Recall:} say the Italian for the cheque, then the Italian for the payment, then say \emph{il deposito} again
\end{itemize}

\begin{tcolorbox}[breakable,colback=teal!4,colframe=teal!35!black,title={Before you move on}]
What does il deposito mean? (\textquotedblleft{}The deposit; the storeroom\textquotedblright{}.) Say it once more.
\end{tcolorbox}
